% GENERATED FILE. Edit canonical lessons, then run npm run generate:books.
% canonical-source-hash: fnv1a64:ad4338dc291f671d
% canonical-lessons: TE-C124-jinka, TE-C124-nakka, TE-C124-pamu, TE-C124-tabelu, TE-C124-cilaka

\chapter{Deer, Fox, Snake, Tortoise, Parrot}
\label{ch:te-deer-fox-snake-tortoise-parrot}
\hlchaptermodality{124}

\begin{chapteropening}
\textbf{By the end of this chapter:} \emph{I can say a deer, a fox, a snake, a tortoise and a parrot.}

Complete the last lesson of chapter 124: i can say a deer, a fox, a snake, a tortoise and a parrot.
\end{chapteropening}

\section[jiṅka]{\te{జింక} (jiṅka) — a deer}

\label{lesson:TE-C124-jinka}



\begin{quote}
Before the new one: say the Telugu for a monkey, then the Telugu for a rabbit.
\end{quote}

\subsection*{You'll want to know: \te{జింక}}
\textbf{\te{జింక}} — \emph{jiṅka} — \textquotedblleft{}a deer\textquotedblright{}.

\begin{grammarlens}[title={What you've built}]
One more word for this chapter.
\end{grammarlens}

\subsection*{Your turn}
\begin{itemize}
\raggedright
  \item \emph{Say it:} \emph{jiṅka}
  \item \emph{Say it:} \emph{jiṅka}, once more
  \item \emph{Say it:} say \emph{jiṅka}
  \item \emph{Recall:} say the Telugu for a monkey, then the Telugu for a rabbit, then say \emph{jiṅka} again
\end{itemize}

\begin{tcolorbox}[breakable,colback=teal!4,colframe=teal!35!black,title={Before you move on}]
What does \te{జింక} mean? (\textquotedblleft{}A deer\textquotedblright{}.) Say it once more.
\end{tcolorbox}
\section[nakka]{\te{నక్క} (nakka) — a fox}

\label{lesson:TE-C124-nakka}



\begin{quote}
Before the new one: say the Telugu for a rabbit, then the Telugu for a deer.
\end{quote}

\subsection*{You'll want to know: \te{నక్క}}
\textbf{\te{నక్క}} — \emph{nakka} — \textquotedblleft{}a fox\textquotedblright{}.

\begin{grammarlens}[title={What you've built}]
One more word for this chapter.
\end{grammarlens}

\subsection*{Your turn}
\begin{itemize}
\raggedright
  \item \emph{Say it:} \emph{nakka}
  \item \emph{Say it:} \emph{nakka}, once more
  \item \emph{Say it:} say \emph{nakka}
  \item \emph{Recall:} say the Telugu for a rabbit, then the Telugu for a deer, then say \emph{nakka} again
\end{itemize}

\begin{tcolorbox}[breakable,colback=teal!4,colframe=teal!35!black,title={Before you move on}]
What does \te{నక్క} mean? (\textquotedblleft{}A fox\textquotedblright{}.) Say it once more.
\end{tcolorbox}
\section[pāmu]{\te{పాము} (pāmu) — a snake}

\label{lesson:TE-C124-pamu}



\begin{quote}
Before the new one: say the Telugu for a deer, then the Telugu for a fox.
\end{quote}

\subsection*{You'll want to know: \te{పాము}}
\textbf{\te{పాము}} — \emph{pāmu} — \textquotedblleft{}a snake\textquotedblright{}.

\begin{grammarlens}[title={What you've built}]
One more word for this chapter.
\end{grammarlens}

\subsection*{Your turn}
\begin{itemize}
\raggedright
  \item \emph{Say it:} \emph{pāmu}
  \item \emph{Say it:} \emph{pāmu}, once more
  \item \emph{Say it:} say \emph{pāmu}
  \item \emph{Recall:} say the Telugu for a deer, then the Telugu for a fox, then say \emph{pāmu} again
\end{itemize}

\begin{tcolorbox}[breakable,colback=teal!4,colframe=teal!35!black,title={Before you move on}]
What does \te{పాము} mean? (\textquotedblleft{}A snake\textquotedblright{}.) Say it once more.
\end{tcolorbox}
\section[tābēlu]{\te{తాబేలు} (tābēlu) — a tortoise}

\label{lesson:TE-C124-tabelu}



\begin{quote}
Before the new one: say the Telugu for a fox, then the Telugu for a snake.
\end{quote}

\subsection*{You'll want to know: \te{తాబేలు}}
\textbf{\te{తాబేలు}} — \emph{tābēlu} — \textquotedblleft{}a tortoise\textquotedblright{}.

\begin{grammarlens}[title={What you've built}]
One more word for this chapter.
\end{grammarlens}

\subsection*{Your turn}
\begin{itemize}
\raggedright
  \item \emph{Say it:} \emph{tābēlu}
  \item \emph{Say it:} \emph{tābēlu}, once more
  \item \emph{Say it:} say \emph{tābēlu}
  \item \emph{Recall:} say the Telugu for a fox, then the Telugu for a snake, then say \emph{tābēlu} again
\end{itemize}

\begin{tcolorbox}[breakable,colback=teal!4,colframe=teal!35!black,title={Before you move on}]
What does \te{తాబేలు} mean? (\textquotedblleft{}A tortoise\textquotedblright{}.) Say it once more.
\end{tcolorbox}
\section[cilaka]{\te{చిలక} (cilaka) — a parrot}

\label{lesson:TE-C124-cilaka}



\begin{quote}
Before the new one: say the Telugu for a snake, then the Telugu for a tortoise.
\end{quote}

\subsection*{You'll want to know: \te{చిలక}}
\textbf{\te{చిలక}} — \emph{cilaka} — \textquotedblleft{}a parrot\textquotedblright{}.

\begin{grammarlens}[title={What you've built}]
One more word for this chapter.
\end{grammarlens}

\subsection*{Your turn}
\begin{itemize}
\raggedright
  \item \emph{Say it:} \emph{cilaka}
  \item \emph{Say it:} \emph{cilaka}, once more
  \item \emph{Say it:} say \emph{cilaka}
  \item \emph{Recall:} say the Telugu for a snake, then the Telugu for a tortoise, then say \emph{cilaka} again
\end{itemize}

\begin{tcolorbox}[breakable,colback=teal!4,colframe=teal!35!black,title={Before you move on}]
What does \te{చిలక} mean? (\textquotedblleft{}A parrot\textquotedblright{}.) Say it once more.
\end{tcolorbox}
